%!TeX root=../main.tex

\chapter{Solução Proposta}

\label{cap:solucao}

Justificada a escolha do OpenTelemetry, falta mostrar o caminho que o usa:
levar o que o kw
registra na máquina de quem o usa até um repositório central, só com
consentimento, sem que o kw deixe de ser um conjunto de comandos curtos em
Bash e sem que cada dado novo obrigue a refazer o caminho. Este capítulo diz o
que cada componente faz e por quê, do lado do cliente ao lado do servidor, e
parte, depois de uma visão geral, do banco do kw, onde o dado nasce, e termina
na extensibilidade do sistema. Como os componentes foram construídos fica para o Capítulo~\ref{cap:implementacao}, e a
verificação, para o Capítulo~\ref{cap:validacao}.

\section{Visão geral da arquitetura}
\label{sec:arquitetura}

O dado percorre quatro componentes até poder ser consultado
(Figura~\ref{fig:arquitetura-caminho}). O kw grava cada evento no seu banco
local, como já fazia (Seções~\ref{sec:funcionalidades}
e~\ref{sec:banco-local}). Um programa separado, o kw\_tele, lê o que ainda não
foi enviado e o manda por OTLP sobre HTTP, com corpo JSON, a um Collector. O
Collector repassa os registros a um backend próprio, escrito para este
trabalho, que os grava num banco SQLite e só então confirma o recebimento; a
confirmação volta pelo mesmo caminho até o kw\_tele, como mostra a
Seção~\ref{sec:backend}. O mesmo backend próprio responde às consultas, lendo desse
banco (Seção~\ref{sec:consulta}).

\begin{figure}[!ht]
  \centering
  \resizebox{\textwidth}{!}{\begin{tikzpicture}[
    x=1cm, y=1cm,
    every node/.style={font=\footnotesize, align=center},
    peca/.style={draw, rounded corners=3pt, minimum width=2.0cm,
                 minimum height=0.85cm, inner sep=2pt, fill=white},
    dado/.style={peca, very thick, fill=black!12},
    lado/.style={draw, dashed, rounded corners=6pt, inner sep=8pt},
    rotulo/.style={font=\scriptsize\scshape, anchor=south west},
    seta/.style={-{Latex[length=1.6mm]}, semithick},
    volta/.style={-{Latex[length=1.6mm]}, semithick, densely dotted},
    nota/.style={font=\scriptsize, inner sep=1pt},
  ]
  \node[peca] (kw) at (0, 2.5) {kw\\[-1pt]{\scriptsize(Bash)}};
  \node[dado] (db) at (0, 0.6) {base local\\[-1pt]{\scriptsize do kw}};
  \node[peca] (tele) at (3.4, 0.6) {kw\_tele\\[-1pt]{\scriptsize(Python)}};
  \node[peca] (cfg) at (3.4, 2.5) {configuração\\[-1pt]{\scriptsize do usuário}};

  \node[peca] (col) at (7.7, 0.6) {Collector};
  \node[peca, minimum width=1.8cm] (ent) at (10.2, 0.6) {entrada};
  \node[peca, minimum width=1.8cm] (cons) at (10.2, 2.5) {consulta};
  \node[draw, very thick, rounded corners=4pt, inner sep=4pt, fit=(ent)(cons)] (back) {};
  \node[nota, font=\scriptsize\bfseries, anchor=north] at ([yshift=-1pt]back.south) {backend próprio};
  \node[dado] (sql) at (12.9, 1.55) {banco\\[-1pt]{\scriptsize SQLite}};
  \node[peca, minimum width=1.7cm] (pesq) at (10.2, 4.1) {usuário};

  \node[lado, fit=(kw)(db)(tele)(cfg)] (maq) {};
  \coordinate (fundo) at ([yshift=-0.45cm]back.south);
  \node[lado, fit=(col)(back)(sql)(pesq)(fundo)] (rec) {};
  \node[rotulo] at ([yshift=2pt]maq.north west) {cliente};
  \node[rotulo] at ([yshift=2pt]rec.north west) {servidor};

  \path (maq.east) -- coordinate[midway] (fx) (rec.west);
  \draw[very thick] ([yshift=4pt]fx |- maq.north) -- ([yshift=-4pt]fx |- maq.south)
    node[below, nota] {fronteira de confiança};

  \draw[seta] (kw) -- node[nota, left] {grava} (db);
  \draw[seta] (tele) -- node[nota, below] {lê e marca} (db);
  \draw[seta] (tele) -- node[nota, right] {lê} (cfg);

  \draw[seta] ([yshift=5pt]tele.east) -- ([yshift=5pt]col.west);
  \draw[seta] ([yshift=5pt]col.east) -- ([yshift=5pt]ent.west);
  \draw[volta] ([yshift=-5pt]ent.west) -- ([yshift=-5pt]col.east);
  \draw[volta] ([yshift=-5pt]col.west) -- ([yshift=-5pt]tele.east);

  \draw[seta] (ent.east -| back.east) -- node[nota, below right, inner sep=0.5pt] {grava} (sql.south west);
  \draw[seta] (cons.east -| back.east) -- node[nota, above right, inner sep=0.5pt] {lê} (sql.north west);
  \draw[seta] (pesq.south) -- node[nota, right] {consulta} (pesq.south |- back.north);
\end{tikzpicture}
}
  \caption{Caminho dos dados, da máquina de quem usa o kw até a consulta. A
    linha vertical marca a fronteira de confiança; a seta pontilhada é a
    confirmação, que volta pelo mesmo caminho depois que o backend próprio
    grava os registros.}
  \label{fig:arquitetura-caminho}
\end{figure}

Do lado do cliente ficam dois componentes. O próprio kw registra as
execuções e, antes de cada comando, decide se pergunta pelo consentimento ou
se dispara um envio em segundo plano, que nunca faz o comando esperar. O
kw\_tele, escrito em Python só com a biblioteca padrão, lê o banco do kw,
reduz a granularidade do tempo, tira os campos que o usuário escolheu não
enviar, monta o registro OTLP, envia e marca o que foi entregue. A
Seção~\ref{sec:cliente} trata dos dois.

Do lado do servidor ficam os outros. O Collector é o receptor de OTLP já
pronto, só configurado: limita o uso de memória e repassa os registros ao
backend, sem filtrar por atributo. O backend próprio separa os registros, põe à
parte os malformados e grava os demais, sem repetir o que já tinha. O backend
também responde às consultas, por rotas só de leitura: umas dizem que tipos de
registro e que atributos existem nos dados, e outra agrupa os registros por
qualquer um desses atributos. Os dois componentes estão na
Seção~\ref{sec:backend}, e a consulta, na Seção~\ref{sec:consulta}.

A fronteira de confiança, a linha entre o que o usuário controla e o que
outros recebem, é a máquina de quem usa o kw. Do lado de dentro estão o banco
do kw, a configuração e o kw\_tele, que roda como o próprio usuário. O único
ponto em que o dado cruza a fronteira é o envio do kw\_tele ao Collector; dali em diante, tudo fica do lado do servidor, e nenhum componente ali sabe o
que cada usuário escolheu.

Por isso, a privacidade é aplicada no cliente, antes de o dado sair da
máquina: o opt-in, a exclusão dos campos que o usuário escolher, a granularidade
de semana do tempo e o identificador aleatório de cada instalação. O Collector
não filtra nada por campo e supõe que o que chega foi consentido. O backend
reduz de novo a granularidade do tempo no que recebe, para que a regra valha
também contra um cliente antigo ou de outra origem. Os detalhes de cada
mecanismo estão no Capítulo~\ref{cap:privacidade}.

As duas pontas têm bancos separados. O do kw segue o esquema do kw; o do
backend é genérico e não sabe nada sobre o kw, porque guarda os atributos que
chegam sem conhecê-los de antemão. Essa separação é o que sustenta a
extensibilidade do sistema, discutida na Seção~\ref{sec:extensibilidade}.

\section{O banco local do kw}
\label{sec:banco-local}

Todo dado que sai da máquina começa no banco do kw
(Seção~\ref{sec:base-local}). O kw grava cada uso como um evento, uma linha
da tabela de eventos com identificador, data e hora, e o resto do que se sabe
sobre aquele uso fica em outras tabelas, ligadas ao evento. Este texto chama
de \emph{tabela de extensão} uma tabela cuja chave aponta, como chave
estrangeira, para o identificador do evento, e que acrescenta informação a
ele; a de execuções é a que o kw já tinha. Os nomes de comando e de estado não
ficam por extenso nessas tabelas: ficam como um número que aponta para uma
\emph{tabela de nomes}. O banco tem um dono só: o kw cria e muda o esquema, e
o kw\_tele apenas lê os eventos e atualiza neles uma marca de envio.

As três propriedades do esquema apresentadas na Seção~\ref{sec:base-local}
não atrapalham o \texttt{kw report}, mas pesam quando o dado sai da máquina,
e a elas se soma uma quarta: o banco não tinha onde registrar que um evento já
foi enviado. Um identificador reaproveitado faria um evento novo passar por um
já enviado. A hora local sem fuso não é o instante absoluto que o OTLP espera.
Sem versão do esquema, nada garante que uma instalação existente tenha uma
coluna nova. E, sem estado de exportação, não havia como evitar reenviar e
contar duas vezes o que já tinha saído. A Tabela~\ref{tab:esquema} resume o
efeito de cada uma e o que foi feito.

\begin{table}[!ht]
  \centering
  \footnotesize
  \setlength{\tabcolsep}{4pt}
  \begin{threeparttable}
    \caption{O que cada propriedade do esquema do kw faz com a exportação e o
      que foi feito.}
    \label{tab:esquema}
    \begin{tabular}{@{}>{\raggedright\arraybackslash}p{4.0cm}
                      >{\raggedright\arraybackslash}p{3.8cm}
                      >{\raggedright\arraybackslash}p{6.4cm}@{}}
      \toprule
      Propriedade do esquema & Efeito na exportação & O que foi feito \\
      \midrule
      identificador por ordem de inserção, sem garantia de não reaproveitamento
        & um evento novo poderia passar por um já enviado
        & contornado: a marca confere identificador e conteúdo, e a chave de
          repetição usa o identificador. Proposta: impedir o reaproveitamento,
          o que exige recriar a tabela \\
      \addlinespace
      data e hora em texto, na hora local, sem fuso
        & a virada da semana fica incerta
        & contornado pela granularidade de semana. Proposta: uma coluna com o
          instante em UTC \\
      \addlinespace
      sem versão do esquema
        & uma coluna nova depende de verificação no ponto de uso
        & feito o mínimo: uma verificação antes do envio, que só acrescenta o que
          falta. Proposta: versão do esquema e uma guarda de migração \\
      \addlinespace
      sem estado de exportação
        & reenvio e contagem dupla
        & feito: uma coluna no evento, com um índice, marca o que foi enviado \\
      \bottomrule
    \end{tabular}
    \begin{tablenotes}
      \footnotesize
      \item Justificativas na Seção~\ref{sec:banco-local}.
    \end{tablenotes}
  \end{threeparttable}
\end{table}


A mudança nas tabelas que já existiam é uma só, e aditiva: uma coluna no evento
que diz se ele está pendente, se ficou de fora do envio por ser anterior ao
consentimento ou em que momento foi enviado. Ela mora no banco do kw, junto com
os dados de uso que descreve, e não num arquivo à parte, e é a própria fila de
envio do kw\_tele (Seção~\ref{sec:cliente}). Como o esquema não tem versão, uma
verificação feita antes de perguntar o consentimento, antes de cada envio e nos
comandos de telemetria acrescenta a coluna aos bancos já instalados quando ela
falta. Nenhuma tabela, coluna ou \textit{view} existente foi alterada. A lista
de módulos entrou como uma tabela de extensão nova, e o kw grava o evento e a
sua linha de extensão juntos, para que a linha nunca fique sem evento
(Seção~\ref{sec:extensibilidade}). Essas mudanças, com o disparador de envio
(Seção~\ref{sec:cliente}), estão num ramo do \textit{fork} do kw usado neste
trabalho e não foram propostas ao projeto.

Três mudanças ficam como proposta. A primeira é uma coluna com o instante do
evento em UTC, ao lado da data e da hora atuais, o que resolveria a virada da
semana e a ordenação, e é defensável como correção isolada. A segunda é
impedir o reaproveitamento do identificador, o que exige recriar a tabela de
eventos. A terceira é versionar o esquema, com uma guarda de migração. O
trabalho parou na mudança mínima porque, sem versionamento, alterar a tabela
de eventos exigiria antes construir a infraestrutura de migração e mexeria em
funções do kw que dependem das \textit{views} atuais, o que, na leitura deste
trabalho, não se justifica só pela telemetria. Enquanto isso, o cliente e o
backend contornam as propriedades: a marca de envio confere o identificador, a
data e a hora do evento, a chave de repetição do backend usa o identificador,
e a granularidade de semana torna o fuso relevante só perto da virada da
semana.

\section{O cliente}
\label{sec:cliente}

O cliente tem duas metades. Uma é pequena, fica dentro do kw, é escrita em
Bash e roda antes de cada comando só para decidir se é preciso fazer alguma
coisa. A outra é um programa separado, o kw\_tele, escrito em Python, que lê
os eventos, prepara o registro e o envia. A divisão segue de uma regra
simples: o que roda em todo comando tem de custar quase nada e nunca falhar, e
montar JSON e falar HTTP com confirmação de entrega é trabalho que o Bash não
faz bem. O kw\_tele é chamado pelo kw do mesmo modo que outras ferramentas
externas que ele já usa, a exemplo do \texttt{patch-hub}
(Seção~\ref{sec:kw-o-que-e}).

Do lado do kw, o registro de uso continua como era: cada execução vira um
evento no banco local, por um único ponto de escrita
(Seção~\ref{sec:funcionalidades}). O que muda é que cada evento ganha uma marca de envio
(Seção~\ref{sec:banco-local}) e que uma verificação curta, o disparador de
envio, roda antes de cada comando. O
disparador faz uma sequência de perguntas (Figura~\ref{fig:disparador}): se o
comando é o próprio \texttt{kw telemetry}, que fica de fora; se o kw\_tele
está instalado; qual é o consentimento; e, com o consentimento dado, se o
comando é dos que registram uso e se já é hora de enviar. Enquanto o usuário
não responde, a pergunta de consentimento vem antes de qualquer outro comando,
inclusive os de ajuda e de versão, que não registram uso.

O disparador decide que é hora de enviar quando venceu o intervalo desde o
último envio bem-sucedido ou quando há eventos pendentes
demais, para que quem usa muito o kw não acumule uma fila grande; entre um
disparo e outro há também um intervalo de tempo mínimo, que cresce quando os
envios falham. O Go faz algo parecido, ao enviar seus relatórios uma vez por semana
(Seção~\ref{sec:precedentes}).

\begin{figure}[!ht]
  \centering
  \begin{tikzpicture}[
    x=1cm, y=1cm,
    every node/.style={font=\footnotesize, align=center},
    passo/.style={draw, rounded corners=3pt, minimum height=0.5cm,
                  inner sep=1.5pt, fill=white, text width=4.0cm},
    pergunta/.style={draw, minimum height=0.5cm, inner sep=1.5pt,
                  fill=black!8, text width=4.4cm},
    fim/.style={passo, very thick},
    seta/.style={-{Latex[length=1.6mm]}, semithick},
    nota/.style={font=\scriptsize, inner sep=1.5pt, fill=white},
  ]
  \node[passo] (ini) at (0, 0) {\texttt{kw} \textit{comando}};
  \node[pergunta] (d1) at (0, -0.8) {o comando é \texttt{telemetry}?};
  \node[pergunta] (d2) at (0, -1.6) {o kw\_tele está instalado?};
  \node[pergunta] (d3) at (0, -2.4) {consentimento?};
  \node[pergunta] (d4) at (0, -3.2) {o comando registra uso?};
  \node[pergunta] (d5) at (0, -4.0) {passou o tempo mínimo?};
  \node[pergunta] (d6) at (0, -4.8) {venceu o intervalo de envio?};
  \node[pergunta] (d7) at (0, -5.6) {há um lote cheio pendente?};
  \node[fim] (segue) at (0, -6.9) {segue para o comando};

  \node[passo, text width=4.6cm] (ask) at (-5.7, -2.4) {pede o consentimento};
  \node[passo, text width=4.6cm] (marca) at (-5.7, -5.6) {guarda o instante do disparo};
  \node[passo, text width=4.6cm] (lanca) at (-5.7, -6.35) {dispara o kw\_tele};

  \draw[seta] (ini) -- (d1);
  \draw[seta] (d1) -- node[nota, right=2pt] {não} (d2);
  \draw[seta] (d2) -- node[nota, right=2pt] {sim} (d3);
  \draw[seta] (d3) -- node[nota, right=2pt] {sim} (d4);
  \draw[seta] (d4) -- node[nota, right=2pt] {sim} (d5);
  \draw[seta] (d5) -- node[nota, right=2pt] {sim} (d6);
  \draw[seta] (d6) -- node[nota, right=2pt] {não} (d7);

  \def\trilho{3.2}
  \draw[semithick] (d1.east) -- node[nota, above=1pt] {sim} (d1.east -| {(\trilho,0)});
  \foreach \d in {d2, d3, d4, d5, d7} {
    \draw[semithick] (\d.east) -- node[nota, above=1pt] {não} (\d.east -| {(\trilho,0)});
  }
  \draw[seta] (d1.east -| {(\trilho,0)}) -- (segue.east -| {(\trilho,0)}) -- (segue.east);

  \draw[seta] (d3.west) -- node[nota, above=1pt, pos=0.62] {perguntar} (ask.east);
  \draw[seta] (ask.west) -- ++(-0.35, 0) |- ([yshift=-3pt]segue.west);

  \draw[seta] (d6.west) -- node[nota, above=1pt] {sim} ([xshift=0.9cm]d6.west -| marca.north) -- ([xshift=0.9cm]marca.north);
  \draw[seta] (d7.west) -- node[nota, above=1pt] {sim} (marca.east);
  \draw[seta] (marca) -- (lanca);
  \draw[seta] (lanca.south) |- ([yshift=3pt]segue.west);
\end{tikzpicture}

  \caption{O que o disparador decide antes de cada comando do kw. Toda saída
    termina no comando; só a pergunta do consentimento espera o usuário, e o
    envio corre em segundo plano.}
  \label{fig:disparador}
\end{figure}

O consentimento é opt-in: o padrão é não enviar nada. O usuário responde a
uma pergunta no terminal, na qual também decide se os eventos gravados antes
da resposta vão junto, pode excluir campos do envio e pode mudar de ideia a
qualquer momento. Tudo isso é aplicado no cliente
(Seção~\ref{sec:arquitetura}), e as respostas possíveis e o que cada uma faz
estão no Capítulo~\ref{cap:privacidade}. O único momento em que o disparador espera o usuário é a pergunta de
consentimento, e só quando há um terminal para respondê-la; as demais
verificações são consultas locais curtas; em todos os casos, ele termina deixando o comando
seguir, e um problema na telemetria nunca impede um comando do kw de rodar.

Quando é hora de enviar, o disparador não envia. Ele dispara o kw\_tele em
segundo plano, com prazo, desligado do terminal, e deixa o comando seguir; o
comando não espera pelo envio. Um backend parado não pode segurar, por
exemplo, uma compilação do kernel. Uma falha no envio vira um estado guardado
na máquina, e não um erro do comando, e o disparador o lê na vez seguinte.
Como não há processo residente, também não há relógio: o envio acontece no
primeiro comando elegível depois de vencido o intervalo, e não numa hora
marcada.

Uma execução do kw\_tele segue um fluxo linear
(Figura~\ref{fig:kwtele-execucao}). Ele confere o consentimento a cada lote,
lê o esquema do banco do kw para saber que tabelas estendem o evento e
seleciona os eventos pendentes em lotes. Para cada evento, reduz a
granularidade do tempo, tira os campos que o usuário excluiu e monta o
registro OTLP (Seção~\ref{sec:mapeamento}); então envia o lote e marca as
linhas só depois de o backend confirmar que as gravou
(Seção~\ref{sec:backend}). A fila de envio é a própria marca no banco do kw, e
não uma fila na memória do programa. Se o backend recusa um lote pelo
conteúdo, o kw\_tele divide o lote até isolar a linha recusada, que fica
pendente para a próxima execução, e segue com as demais; ao fim, a execução
conta como falha, para que a linha seja tentada de novo; qualquer outra falha
interrompe a execução, que é refeita numa vez seguinte. O programa não conhece
as colunas de antemão: ele as lê do banco, o que é a base do argumento da
Seção~\ref{sec:extensibilidade}.

\begin{figure}[!ht]
  \centering
  \resizebox{\textwidth}{!}{\begin{tikzpicture}[
    x=1cm, y=1cm,
    every node/.style={font=\footnotesize, align=center},
    passo/.style={draw, rounded corners=3pt, minimum height=0.7cm,
                  inner sep=3pt, fill=white, text width=4.0cm},
    pergunta/.style={draw, minimum height=0.7cm, inner sep=3pt,
                  fill=black!8, text width=4.0cm},
    fim/.style={passo, very thick, text width=3.9cm},
    lado/.style={draw, dashed, rounded corners=6pt, inner sep=7pt},
    rotulo/.style={font=\scriptsize\scshape, anchor=south west},
    seta/.style={-{Latex[length=1.6mm]}, semithick},
    rede/.style={-{Latex[length=1.6mm]}, semithick, densely dotted},
    nota/.style={font=\scriptsize, inner sep=1.5pt},
  ]
  \node[passo] (esq) at (0, 0) {lê o esquema da base local};
  \node[passo] (pend) at (0, -1.3) {lê um lote de\\eventos pendentes};
  \node[passo] (tempo) at (0, -2.6) {reduz a granularidade do tempo à semana};
  \node[passo] (excl) at (0, -3.9) {exclui os campos da lista do usuário};
  \node[passo] (monta) at (0, -5.2) {monta o registro OTLP};
  \node[passo] (envia) at (0, -6.5) {envia o lote};
  \node[pergunta] (resp) at (0, -7.8) {resposta?};
  \node[passo] (marca) at (0, -9.4) {marca o lote como enviado};

  \node[fim] (termina) at (7.0, -1.3) {guarda o estado e termina};
  \node[fim] (para) at (7.0, -9.4) {para e guarda o estado};
  \node[passo, text width=3.4cm] (divide) at (-6.0, -7.8)
    {divide o lote até isolar a linha recusada, que fica pendente};

  \coordinate (vl) at (-8.5, -4.5);
  \node[lado, fit=(esq)(marca)(termina)(para)(divide)(vl)] (maq) {};
  \node[rotulo] at ([yshift=2pt]maq.north west) {na máquina de quem usa o kw};

  \node[passo, text width=2.4cm, minimum height=2.0cm] (rec) at (12.0, -7.15)
    {Collector e\\backend próprio};
  \node[rotulo] at ([yshift=2pt]rec.north west) {pela rede};

  \draw[seta] (esq) -- (pend);
  \draw[seta] (pend) -- (tempo);
  \draw[seta] (tempo) -- (excl);
  \draw[seta] (excl) -- (monta);
  \draw[seta] (monta) -- (envia);
  \draw[seta] (envia) -- (resp);
  \draw[seta] (pend) -- node[nota, above] {nenhum pendente} (termina);

  \draw[rede] (envia.east) -- (envia.east -| rec.west);
  \draw[rede] ([yshift=5pt]resp.east -| rec.west) -- node[nota, above, pos=0.3] {resposta} ([yshift=5pt]resp.east);

  \draw[seta] (resp) -- node[nota, right] {aceito} (marca);
  \draw[seta] (resp.west) -- node[nota, above, align=center] {recusado\\pelo conteúdo} (divide.east);
  \draw[seta] ([yshift=-5pt]resp.east) -| node[nota, below, pos=0.25] {outra falha} (para.north);

  \def\volta{-8.5}
  \draw[seta] (marca.west) -- (\volta, -9.4) -- (\volta, -1.3) -- (pend.west);
  \draw[semithick] (divide.west) -- (\volta, -7.8);
  \node[nota, rotate=90, below] at (\volta, -4.5) {próximo lote};
\end{tikzpicture}
}
  \caption{Uma execução do kw\_tele. As setas pontilhadas atravessam a rede; um
    lote recusado pelo conteúdo é dividido até isolar a linha recusada, que fica
    pendente.}
  \label{fig:kwtele-execucao}
\end{figure}

Esse design decorre do que o kw é: uma ferramenta em Bash, que não deve ganhar
dependências pesadas, feita de comandos curtos, sem nada residente, e cuja
telemetria não pode atrasar nem quebrar um comando. Disso decorrem duas
escolhas. A primeira é não usar o SDK do OpenTelemetry no cliente. O SDK traz
instrumentação e convenções prontas, mas, no SDK de Python com processamento
em lote, quem envia não sabe se o envio chegou, e este design precisa marcar
cada linha como enviada só depois da confirmação; além disso, o SDK traria
dependências que o cliente evita. A comparação, com as fontes, está nas
Seções~\ref{sec:alternativas} e~\ref{sec:limites-otel}; o caminho síncrono do
SDK não foi avaliado.

A segunda escolha é usar o OTLP sobre HTTP com corpo JSON, e não sobre gRPC
com protobuf. O Python traz JSON e HTTP na biblioteca padrão e não traz gRPC
nem protobuf, e o corpo em JSON pode ser lido por quem quiser conferir o que
seria enviado. O preço é um corpo maior e as regras do JSON do OTLP,
discutidas na Seção~\ref{sec:limites-otel}. Na leitura deste trabalho, um
processo que roda por segundos também não se beneficiaria da conexão longa do
gRPC.

\section{Do evento do kw ao registro OTLP}
\label{sec:mapeamento}

Esta seção acompanha um evento do kw do momento em que ele é gravado até sair
da máquina como registro OTLP. O exemplo é uma compilação bem-sucedida com a
lista de módulos, a fonte acrescentada ao kw como demonstração
(Seção~\ref{sec:extensibilidade}). Quando o usuário compila o kernel, o kw grava
um evento: uma linha da tabela de eventos, com identificador, data e hora
(Seção~\ref{sec:base-local}). O resto do que se sabe sobre aquela compilação fica nas tabelas de extensão
do evento (Seção~\ref{sec:banco-local}). Aqui há duas. A da execução, com o comando, a duração e o estado em que
terminou, já existia no kw. A da lista de módulos, com a contagem e os nomes,
é a nova, acrescentada neste exemplo. O comando e o estado ficam como um número que aponta para uma tabela de nomes.

Na hora do envio, o kw\_tele não sabe de antemão que tabelas existem. Ele lê o
esquema do banco e acha as tabelas de extensão por essa chave estrangeira, de
modo que uma tabela nova ligada ao evento passa a ser enviada sem mudar o
cliente. Para cada evento pendente, junta a linha do evento com as linhas das
extensões que têm o mesmo identificador e troca cada número pelo nome
correspondente. Cada linha de extensão vira um registro OTLP, e um evento sem
extensão vira um registro só. Nem todo comando preenche todas as tabelas de
extensão: cada comando grava só as que lhe dizem respeito. Compilar o kernel ou
instalá-lo numa máquina de teste preenche a da execução e a dos módulos; listar
ou remover os kernels instalados, só a da execução; o \texttt{pomodoro}
preenche apenas a sua, de sessões de trabalho. O número de registros por evento
varia, portanto, com o comando. A compilação do exemplo vira, portanto, dois
registros com o mesmo identificador de evento
(Figura~\ref{fig:evento-registro}).

\begin{figure}[!ht]
  \centering
  \resizebox{\textwidth}{!}{\begin{tikzpicture}[
    x=1cm, y=1cm,
    every node/.style={font=\footnotesize},
    tabela/.style={draw, rectangle split, rectangle split parts=2,
                   rectangle split part fill={black!12, white},
                   inner sep=3pt, font=\footnotesize},
    registro/.style={draw, rectangle split, rectangle split parts=2,
                   rectangle split part fill={black!6, white},
                   inner sep=2pt, font=\footnotesize, rounded corners=2pt},
    recurso/.style={draw, very thick, rounded corners=4pt, inner sep=3pt},
    coluna/.style={draw, dashed, rounded corners=6pt, inner sep=5pt},
    rotulo/.style={font=\scriptsize\scshape, anchor=south west},
    seta/.style={-{Latex[length=1.6mm]}, semithick},
    nota/.style={font=\footnotesize, align=left, inner sep=1.5pt},
  ]
  \providecommand{\kwcolunas}{}\renewcommand{\kwcolunas}[1]{%
    {\setlength{\tabcolsep}{0pt}\begin{tabular}{@{}l<{\hspace{0.45em}}@{}r@{}}#1\end{tabular}}}
  \providecommand{\kwcol}{}\renewcommand{\kwcol}[2]{\texttt{#1} & {\scriptsize #2}\\}
  \providecommand{\kwchave}{}\renewcommand{\kwchave}[2]{\rowcolor{black!12}\texttt{\textbf{#1}} & {\scriptsize\textbf{#2}}\\}

  \node[tabela, anchor=north west] (ev) at (0, 0) {%
    \texttt{event}
    \nodepart{two}
    \kwcolunas{\kwchave{id}{4127}\kwcol{date}{2026-11-19}\kwcol{time}{14:32:07}\kwcol{sent\_at}{---}}};

  \node[tabela, anchor=north west] (ex) at ([yshift=-0.35cm]ev.south west) {%
    \texttt{executed}
    \nodepart{two}
    \kwcolunas{\kwchave{id $\to$ event.id}{4127}\kwcol{command\_label\_id}{3}%
             \kwcol{elapsed\_time\_in\_secs}{418}\kwcol{status\_id}{1}}};

  \node[tabela, anchor=north west] (bm) at ([yshift=-0.35cm]ex.south west) {%
    \texttt{build\_modules}
    \nodepart{two}
    \kwcolunas{\kwchave{id $\to$ event.id}{4127}\kwcol{module\_count}{3}%
             \kwcol{module\_names}{\texttt{amdgpu,\dots}}}};

  \node[nota, anchor=north west] (nomes) at ([yshift=-0.2cm]bm.south west) {nomes: \texttt{3} = \texttt{build},
    \texttt{1} = \texttt{success}};


  \node[coluna, fit=(ev)(ex)(bm)(nomes)] (g) {};
  \node[rotulo] at ([yshift=2pt]g.north west) {geração: o banco do kw};

  \coordinate (tx) at ([xshift=1.0cm]g.east);
  \node[nota, anchor=west, text width=5.2cm] (t1) at (tx |- ev.center)
    {data e hora $\to$ instante da semana; a data vira a segunda-feira; hora e marca de envio não saem};
  \node[nota, anchor=west, text width=5.2cm] (t3) at (tx |- ex.center) {número $\to$ nome, pelas tabelas de nomes};
  \node[nota, anchor=west, text width=5.2cm] (t4) at (tx |- bm.center)
    {acha as tabelas de extensão pela chave para \texttt{event.id}; cada linha de extensão vira um registro};
  \node[nota, anchor=north west, text width=5.2cm] (t5) at ([yshift=-0.25cm]t4.south west)
    {campos da lista do usuário saem do registro};

  \node[coluna, fit=(t1)(t5), inner sep=6pt] (tr) {};
  \node[rotulo] at ([yshift=2pt]tr.north west) {transformação: o kw\_tele};

  \draw[seta] (ev.east -| g.east) -- (t1.west);
  \draw[seta] (ex.east -| g.east) -- (t3.west);
  \draw[seta] (bm.east -| g.east) -- (t4.west);

  \node[registro, anchor=north west] (r1) at ([yshift=-1.1cm]g.south west) {%
    registro 1 \quad corpo: \texttt{executed}
    \nodepart{two}
    \kwcolunas{\kwcol{kw.event.id}{4127}\kwcol{kw.event.date}{2026-11-16}%
             \kwcol{kw.event.kind}{executed}%
             \kwcol{kw.executed.command\_label}{build}%
             \kwcol{kw.executed.elapsed\_time\_in\_secs}{418}\kwcol{kw.executed.status}{success}}};

  \node[registro, anchor=north west] (r2) at ([xshift=0.15cm]r1.north east) {%
    registro 2 \quad corpo: \texttt{build\_modules}
    \nodepart{two}
    \kwcolunas{\kwcol{kw.event.id}{4127}\kwcol{kw.event.date}{2026-11-16}%
             \kwcol{kw.event.kind}{build\_modules}%
             \kwcol{kw.build\_modules.module\_count}{3}%
             \kwcol{kw.build\_modules.module\_names}{\texttt{amdgpu,\dots}}}};

  \node[recurso, fit=(r1)(r2)] (rq) {};
  \node[rotulo] at ([yshift=2pt]rq.north west) {trânsito: os registros, sob um só recurso};

  \draw[seta] (tr.south) -- (tr.south |- rq.north);
\end{tikzpicture}
}
  \caption{Um evento de compilação do kw virando dois registros OTLP, um por
    linha de extensão, sob um só recurso (Figura~\ref{fig:requisicao-otlp}). Valores
    ilustrativos.}
  \label{fig:evento-registro}
\end{figure}

Cada registro é montado com o que foi lido. O corpo traz o tipo do registro,
isto é, o nome da tabela de extensão, ou \enquote{event} quando não há
extensão, e o mesmo valor vai num atributo de tipo, para que a consulta agrupe
por tipo. Cada coluna vira um atributo cujo nome diz de onde ela veio:
\texttt{kw.event.<coluna>} para o evento e \texttt{kw.<tabela>.<coluna>} para a
extensão, como \texttt{kw.build\_modules.module\_count} no registro dos
módulos. A coluna que apontava para uma tabela de nomes sai com o nome, como o
comando, em \texttt{kw.executed.command\_label}. O identificador do evento vai
como atributo, porque o registro OTLP não tem identidade própria e o backend
precisa dele para reconhecer repetidos (Seção~\ref{sec:backend}). O nível de
severidade, o campo do modelo de \textit{logs} que classifica o registro como
informação, aviso ou erro, não tem equivalente no kw e sai com o valor neutro
de informação. A Tabela~\ref{tab:mapeamento} lista todos os campos.

\begin{table}[!ht]
  \centering
  \footnotesize
  \setlength{\tabcolsep}{4pt}
  \begin{threeparttable}
    \caption{Do evento do kw ao registro OTLP. As colunas de extensão não são
      fixas: o cliente as descobre no banco do kw a cada execução, e
      \texttt{executed} e \texttt{build\_modules} são exemplos.}
    \label{tab:mapeamento}
    \begin{tabular}{@{}>{\raggedright\arraybackslash}p{4.0cm}
                      >{\raggedright\arraybackslash}p{5.0cm}
                      >{\raggedright\arraybackslash}p{5.2cm}@{}}
      \toprule
      Origem no kw & Campo no registro OTLP & Tratamento \\
      \midrule
      \multicolumn{3}{@{}l}{\textit{Evento}} \\*
      \addlinespace[1pt]
      identificador do evento & \texttt{kw.event.id} & direto \\
      outra coluna do evento & \texttt{kw.event.}\textit{coluna} & direto \\
      data do evento & \texttt{kw.event.date} & vira a segunda-feira da semana \\
      hora do evento & --- & não enviada \\
      marca de envio & --- & não enviada \\
      data e hora do evento & \texttt{timeUnixNano} & instante da semana \\
      momento do envio & \texttt{observedTimeUnixNano} & instante da semana \\
      \midrule
      \multicolumn{3}{@{}l}{\textit{Extensões (descobertas no banco do kw)}} \\*
      \addlinespace[1pt]
      tabela de extensão & \texttt{body} e \texttt{kw.event.kind}
        & nome da tabela; \texttt{event} quando o evento não tem extensão;
          um registro por linha de extensão \\
      coluna da extensão & \texttt{kw.}\textit{tabela}\texttt{.}\textit{coluna} & direto \\
      coluna que aponta para uma tabela de nomes
        & \texttt{kw.}\textit{tabela}\texttt{.}\textit{coluna} sem o sufixo \texttt{\_id}
          (por exemplo, \texttt{kw.executed.command\_label})
        & o nome, no lugar do número \\
      \midrule
      \multicolumn{3}{@{}l}{\textit{Recurso (um por envio)}} \\*
      \addlinespace[1pt]
      instalação & \texttt{service.instance.id} & UUID aleatório, criado no opt-in \\
      kw e cliente
        & \texttt{service.name}, \texttt{service.version},
          \texttt{telemetry.sdk.name}, \texttt{.language}, \texttt{.version}
        & fixos, exceto as versões \\
      \midrule
      sem equivalente no kw & \texttt{severityNumber}, \texttt{severityText}
        & fixos: 9 e \texttt{INFO} \\
      \bottomrule
    \end{tabular}
    \begin{tablenotes}
      \footnotesize
      \item Justificativas na Seção~\ref{sec:mapeamento}.
    \end{tablenotes}
  \end{threeparttable}
\end{table}


O tempo é o único campo que muda de valor no caminho, e muda de propósito. No
kw, a data e a hora do evento estão em texto, no horário local e sem fuso. No
registro, o tempo vai com granularidade de semana: o campo do tempo do evento
leva o instante da semana em que ele ocorreu, um mesmo instante para toda a
semana, e o atributo da data passa a ser a segunda-feira daquela semana. A
hora não sai, e a marca de envio também não. O motivo é que o horário de
trabalho do usuário não sai da máquina (Capítulo~\ref{cap:privacidade}). A
redução não é configurável, e, quando a data guardada não é legível, o tempo
sai como desconhecido e a data não é enviada. O modelo de \textit{logs} prevê
nesse campo o instante em que o evento ocorreu \citep{otel_logs_data_model};
levar o instante da semana é um desvio deliberado, e um evento perto da virada
da semana em UTC pode cair na semana vizinha, limite retomado na
Seção~\ref{sec:banco-local}.

Os registros de um envio saem juntos, sob um recurso, a parte do documento OTLP
que diz quem os produziu: o serviço, a versão do kw, um identificador da
instalação e o nome, a linguagem e a versão do cliente, com os nomes de
atributo das convenções do OpenTelemetry, que marcam os quatro últimos como
estáveis \citep{otel_semconv_resource_v1_45_0}. O identificador da instalação é
aleatório, criado quando o usuário aceita enviar, como as convenções
recomendam \citep{otel_semconv_registry_service_v1_45_0}; ele serve para contar
instalações e separar os repetidos de cada uma, sem relação com o usuário ou a
máquina (Capítulo~\ref{cap:privacidade}). O cliente se identifica como
kw\_tele, e não pelo nome que as convenções reservam ao SDK oficial
\citep{otel_semconv_resource_v1_45_0}. Cada registro leva também o momento do
envio, como momento da observação, igualmente com granularidade de semana. O documento segue por HTTP, com corpo JSON, para o Collector
(Figura~\ref{fig:requisicao-otlp}). Cada evento vira um
registro de \textit{log} porque é um fato com tempo e atributos; contar e
agrupar fica para o backend.

\begin{figure}[!ht]
  \centering
  \begin{tikzpicture}[
    x=1cm, y=1cm,
    every node/.style={font=\footnotesize},
    linha/.style={anchor=west, inner sep=1pt, font=\footnotesize\ttfamily},
    titulo/.style={anchor=north west, inner sep=2pt, xshift=2pt, yshift=-1pt, font=\footnotesize\ttfamily\bfseries},
    nivel/.style={draw, rounded corners=3pt, inner sep=5pt},
    recolhido/.style={anchor=west, inner sep=1pt, font=\footnotesize\itshape, text=black!65},
    nota/.style={anchor=west, font=\scriptsize, align=left, inner sep=1.5pt, text width=4.8cm},
    liga/.style={semithick, densely dotted},
  ]
  \providecommand{\aspas}{}\renewcommand{\aspas}[1]{\textquotedbl{}#1\textquotedbl{}}

  \node[linha] (http) at (0, 0) {POST \textrm{\textit{<endpoint configurado>}}};
  \node[linha] (ct) at (0, -0.45) {Content-Type: application/json};

  \node[linha] (a1) at (0.9, -2.0) {service.name \textrm{=} \aspas{kworkflow}};
  \node[linha] (a2) at (0.9, -2.4) {service.version \textrm{=} \textrm{\textit{<versão do kw>}}};
  \node[linha] (a3) at (0.9, -2.8) {service.instance.id \textrm{=} \textrm{\textit{<UUID aleatório>}}};
  \node[linha] (a4) at (0.9, -3.2) {telemetry.sdk.name \textrm{=} \aspas{kw\_tele}};
  \node[linha] (a5) at (0.9, -3.6) {telemetry.sdk.language \textrm{=} \aspas{python}};
  \node[linha] (a6) at (0.9, -4.0) {telemetry.sdk.version \textrm{=} \textrm{\textit{<versão do kw\_tele>}}};
  \node[nivel, fit=(a1)(a6), inner ysep=9pt, yshift=6pt] (fres) {};
  \node[titulo] at (fres.north west) {resource.attributes};

  \node[linha] (sc) at (1.3, -5.25) {scope \textrm{=} \{name: \aspas{kw\_tele}, version\}};
  \node[linha] (r1) at (1.7, -6.25) {timeUnixNano, observedTimeUnixNano};
  \node[linha] (r2) at (1.7, -6.65) {severityNumber \textrm{=} 9};
  \node[linha] (r2b) at (1.7, -7.05) {severityText \textrm{=} \aspas{INFO}};
  \node[linha] (r3) at (1.7, -7.45) {body.stringValue \textrm{=} \textrm{\textit{<tipo>}}};
  \node[recolhido] (r4) at (1.7, -7.85) {attributes: \texttt{kw.event.*}, \texttt{kw.}\textit{<tabela>}\texttt{.*}};
  \node[recolhido] (r4b) at (1.7, -8.2) {(Figura~\ref{fig:evento-registro})};
  \node[nivel, fit=(r1)(r4b), inner ysep=9pt, yshift=6pt] (frec) {};
  \node[titulo] at (frec.north west) {logRecords[0]};
  \node[recolhido] (r5) at (1.3, -9.0) {logRecords[1], \dots};
  \node[nivel, fit=(sc)(frec)(r5), inner ysep=9pt, yshift=6pt] (fsl) {};
  \node[titulo] at (fsl.north west) {scopeLogs[0]};

  \node[nivel, very thick, fit=(fres)(fsl), inner ysep=10pt, yshift=7pt] (frl) {};
  \node[titulo] at (frl.north west) {resourceLogs[0]};

  \coordinate (col) at ([xshift=0.6cm]frl.east);
  \node[nota] (n0) at ([yshift=-0.2cm]col |- http) {a requisição: JSON, sem compressão,\\ao endpoint do Collector\\(o padrão termina em \texttt{/v1/logs})};
  \node[nota] (n1) at (col |- a1) {\textbf{resource}: uma vez por requisição; quem produziu os dados};
  \node[nota] (n2) at (col |- a3) {o identificador da instalação: aleatório, criado ao aceitar};
  \node[nota] (n3) at ([yshift=0.15cm]col |- sc) {\textbf{scope}: o cliente que montou\\os registros};
  \node[nota] (n4) at ([yshift=-0.3cm]col |- r1) {\textbf{logRecords}: um por linha\\de extensão; os dois tempos vão na\\granularidade de semana: o do\\evento e o do envio};
  \node[nota] (n5) at ([yshift=-0.75cm]col |- r3) {o corpo é o tipo do registro (o nome da tabela)};

  \draw[liga] (a1.east) -- (n1.west);
  \draw[liga] (a3.east) -- (n2.west);
  \draw[liga] (sc.east) -- (sc.east -| n3.west);
  \draw[liga] (r1.east) -| ([xshift=-0.3cm]n4.west) -- (n4.west);
  \draw[liga] (r3.east) -| ([xshift=-0.3cm]n5.west) -- (n5.west);
\end{tikzpicture}

  \caption{Uma requisição OTLP decomposta. Os atributos de cada registro estão
    na Figura~\ref{fig:evento-registro}.}
  \label{fig:requisicao-otlp}
\end{figure}

Na tradução, o evento virou o registro, cada coluna virou um atributo e a
tabela virou o tipo; pediram adaptação o tempo, as colunas de nomes, o
identificador do evento e o nível de severidade. O que o kw não guarda, o fuso
e a versão do esquema, e o identificador que pode ser reaproveitado são o
assunto da Seção~\ref{sec:banco-local}.

\section{O Collector e o backend próprio}
\label{sec:backend}

O Collector recebe, processa e repassa registros, mas não os guarda num banco
de dados. A distribuição usada, a principal, na versão 0.121.0, não traz
exportador para banco: os seus exportadores enviam a outro receptor, a um
arquivo ou a sistemas de métricas e de mensagens
\citep{otelcol_releases_manifest_0_121_0}. Exportadores para bancos, como o do
ClickHouse, existem só na distribuição \textit{contrib}, em estágio beta para
\textit{logs} \citep{otelcol_contrib_clickhouse_metadata_0_121_0}, e nenhum
grava em SQLite. Por isso este trabalho pôs um backend próprio depois do
Collector, e a divisão de papéis ficou assim: o Collector é a parte exposta
aos clientes, que recebe o protocolo e se protege de picos de carga; o backend
fica atrás dele, separa, guarda e confirma, e serve a consulta
(Seção~\ref{sec:consulta}). Na leitura deste trabalho, é essa divisão que
deixa ao código escrito aqui só o que é específico do caminho.

O Collector é usado como vem, só por configuração
(Figura~\ref{fig:collector}): um receptor de OTLP sobre
HTTP, um limite de uso de memória, um processador que só acrescenta um
atributo de marcação e um exportador que repassa os registros ao backend
próprio. É esse reuso que motivou a escolha do OpenTelemetry
(Seção~\ref{sec:criterios}): o trabalho não precisou escrever o receptor
voltado aos clientes nem a proteção contra picos de carga. O limite de memória
vem antes de tudo e recusa o excesso antes de processá-lo, e o cliente trata
essa recusa como passageira, deixando as linhas para uma próxima vez. O
Collector não agrupa registros em lote nem guarda fila, por uma razão
explicada adiante, e não filtra nada por campo: ele supõe que o que chega já
foi consentido (Seção~\ref{sec:arquitetura}).

\begin{figure}[!ht]
  \centering
  \resizebox{\textwidth}{!}{\begin{tikzpicture}[
    x=1cm, y=1cm,
    every node/.style={font=\footnotesize, align=center},
    peca/.style={draw, rounded corners=3pt, minimum height=0.85cm,
                 inner sep=3pt, fill=white, text width=2.6cm},
    ausente/.style={draw, densely dotted, rounded corners=3pt, minimum height=0.7cm,
                 inner sep=3pt, text width=2.6cm, font=\scriptsize},
    lado/.style={draw, dashed, rounded corners=6pt, inner sep=8pt},
    rotulo/.style={font=\scriptsize\scshape, anchor=south west},
    seta/.style={-{Latex[length=1.6mm]}, semithick},
    volta/.style={-{Latex[length=1.6mm]}, semithick, densely dotted},
    nota/.style={font=\scriptsize, inner sep=1.5pt},
  ]
  \node[peca] (rec) at (0, 0) {receptor\\[-1pt]{\scriptsize OTLP/HTTP}};
  \node[peca] (mem) at (3.6, 0) {limite de memória};
  \node[peca] (exp) at (7.2, 0) {exportador\\[-1pt]{\scriptsize OTLP/HTTP}};

  \node[ausente] (a1) at (1.8, -1.7) {sem lote};
  \node[ausente] (a2) at (5.4, -1.7) {sem fila nem novas tentativas};
  \node[ausente] (a3) at (3.6, 1.6) {sem filtro por campo};

  \node[lado, fit=(rec)(exp)(a1)(a2)(a3)] (col) {};
  \node[rotulo] at ([yshift=2pt]col.north west) {Collector};

  \node[nota, anchor=east] (in) at (-2.6, 0) {envio do\\kw\_tele};
  \node[nota, anchor=west] (out) at (9.8, 0) {envio ao\\backend próprio};

  \draw[seta] (in.east) -- (rec.west);
  \draw[seta] (rec) -- (mem);
  \draw[seta] (mem) -- (exp);
  \draw[seta] (exp.east) -- (out.west);

  \draw[volta] (mem.south) -- ++(0, -0.45) coordinate (m1) -- (m1 -| in.east)
    node[nota, left, align=right] {recusa passageira:\\as linhas ficam\\pendentes};
\end{tikzpicture}
}
  \caption{Dentro do Collector, no caminho dos registros. O que está
    pontilhado não existe neste caminho.}
  \label{fig:collector}
\end{figure}

O backend próprio entende só a pequena parte do OTLP sobre HTTP de que o
caminho precisa, que é receber registros de \textit{log}. Ele separa o corpo da
requisição que o Collector repassa em registros
(Figura~\ref{fig:backend-entrada}), e o que não se encaixa no formato vai para um
armazenamento à parte, junto com o motivo, sem derrubar o resto do lote. O
conteúdo nunca é motivo para recusar uma requisição: se fosse, cada registro
malformado faria o cliente dividir o lote até isolá-lo e deixá-lo pendente,
para recusá-lo de novo a cada tentativa.

\begin{figure}[!ht]
  \centering
  \resizebox{\textwidth}{!}{\begin{tikzpicture}[
    x=1cm, y=1cm,
    every node/.style={font=\footnotesize, align=center},
    passo/.style={draw, rounded corners=3pt, minimum height=0.7cm,
                  inner sep=3pt, fill=white, text width=3.6cm},
    pergunta/.style={draw, minimum height=0.7cm, inner sep=3pt,
                  fill=black!8, text width=3.6cm},
    dado/.style={passo, very thick, fill=black!12},
    fim/.style={passo, very thick},
    lado/.style={draw, dashed, rounded corners=6pt, inner sep=8pt},
    rotulo/.style={font=\scriptsize\scshape, anchor=south west},
    seta/.style={-{Latex[length=1.6mm]}, semithick},
    volta/.style={-{Latex[length=1.6mm]}, semithick, densely dotted},
    nota/.style={font=\scriptsize, inner sep=1.5pt},
  ]
  \node[passo] (corpo) at (0, 0) {corpo da requisição};
  \node[passo] (separa) at (0, -1.2) {separa em registros};
  \node[pergunta] (bom) at (0, -2.4) {o registro está bem formado?};
  \node[passo] (tempo) at (-2.6, -3.8) {reduz a granularidade do tempo à semana};
  \node[passo] (dup) at (-2.6, -5.1) {calcula a chave de repetição};
  \node[passo] (aparte) at (2.6, -4.45) {guarda à parte, com o motivo};
  \node[dado, text width=4.4cm] (grava) at (0, -6.6) {grava o lote de uma vez, ignorando os repetidos};
  \node[dado, text width=1.9cm] (sql) at (4.6, -6.6) {banco\\[-1pt]{\scriptsize SQLite}};
  \node[fim] (ok) at (-2.6, -8.4) {responde sucesso};
  \node[fim] (erro) at (2.6, -8.4) {responde erro};

  \node[lado, fit=(corpo)(tempo)(aparte)(sql)(ok)(erro)] (back) {};
  \node[rotulo] at ([yshift=2pt]back.north west) {backend próprio};

  \draw[seta] (corpo) -- (separa);
  \draw[seta] (separa) -- (bom);
  \draw[seta] (bom.west) -| node[nota, above, pos=0.25] {sim} (tempo.north);
  \draw[seta] (bom.east) -| node[nota, above, pos=0.25] {não} (aparte.north);
  \draw[seta] (tempo) -- (dup);
  \draw[seta] (dup.south) -- ++(0, -0.45) -| ([xshift=-0.6cm]grava.north);
  \draw[seta] (aparte.south) -- ++(0, -1.1) -| ([xshift=0.6cm]grava.north);
  \draw[seta] (grava) -- (sql);
  \draw[seta] (grava.south) -- ++(0, -0.35) -| node[nota, left, pos=0.75] {gravou} (ok.north);
  \draw[seta] (grava.south) -- ++(0, -0.35) -| node[nota, right, pos=0.75] {falha de armazenamento} (erro.north);

  \node[passo, text width=2.4cm] (exp) at (-10.2, 0) {exportador\\[-1pt]{\scriptsize OTLP/HTTP}};
  \node[nota, text width=2.6cm, anchor=north] (repassa) at ([yshift=-0.35cm]exp.south) {a resposta é repassada\\ao kw\_tele};
  \node[lado, fit=(exp)(repassa)] (col) {};
  \node[rotulo] at ([yshift=2pt]col.north west) {Collector};

  \draw[seta] ([yshift=4pt]exp.east) -- node[nota, above, pos=0.24, align=center] {POST \texttt{/v1/logs}\\OTLP, JSON} ([yshift=4pt]corpo.west);

  \coordinate (trilho) at ([xshift=-1.0cm]back.west);
  \draw[volta] (ok.west) -- node[nota, above, pos=0.7] {sucesso} (ok.west -| trilho);
  \draw[volta] (erro.south) -- ([yshift=-0.35cm]erro.south |- back.south) coordinate (e1) -- node[nota, below, pos=0.5] {erro} (e1 -| trilho);
  \draw[volta] (e1 -| trilho) |- ([yshift=-4pt]exp.east);
\end{tikzpicture}
}
  \caption{A entrada do backend próprio. A granularidade do tempo é reduzida à
    semana antes do cálculo da chave de repetição, e o lote é gravado de uma
    vez, sem os repetidos, antes da resposta.}
  \label{fig:backend-entrada}
\end{figure}

O backend próprio guarda os dados em três tabelas
(Figura~\ref{fig:esquema-backend}). Cada requisição recebida vira uma linha numa
tabela de lotes, com o momento da chegada, em granularidade de semana, e o
tamanho do corpo. Cada registro bem formado vira uma linha numa tabela de
registros, ligada ao seu lote: ela tem uma coluna para cada campo fixo do OTLP,
como a instalação, o corpo, o tempo do evento, o da observação e o
identificador do evento, e guarda todos os atributos juntos, num único
documento JSON. Os registros malformados vão para uma terceira tabela, também
ligada ao lote, com o motivo e o corpo como chegou. Não há uma coluna por
atributo nem uma tabela por fonte: um atributo novo é só mais uma chave no
documento, e a consulta o lê pelo nome (Seção~\ref{sec:extensibilidade}). O
preço é que o banco não sabe o tipo de cada atributo: a consulta agrupa pelos
valores lidos como texto e só soma uma medida onde o valor é numérico,
contando à parte os que não são.

\begin{figure}[!ht]
  \centering
  \begin{tikzpicture}[
    x=1cm, y=1cm,
    every node/.style={font=\footnotesize},
    tabela/.style={draw, rectangle split, rectangle split parts=2,
                   rectangle split part fill={black!12, white},
                   inner sep=3pt, font=\footnotesize},
    fk/.style={-{Latex[length=1.6mm]}, semithick},
    nota/.style={font=\scriptsize, align=left, inner sep=1.5pt},
  ]
  \providecommand{\kwcolunas}{}\renewcommand{\kwcolunas}[1]{%
    {\setlength{\tabcolsep}{0pt}\begin{tabular}{@{}l<{\hspace{1.2em}}@{}r@{}}#1\end{tabular}}}
  \providecommand{\kwcol}{}\renewcommand{\kwcol}[2]{\texttt{#1} & {\scriptsize #2}\\}
  \providecommand{\kwdestaque}{}\renewcommand{\kwdestaque}[2]{%
    \rowcolor{black!12}\texttt{\textbf{#1}} & {\scriptsize\textbf{#2}}\\}

  \node[tabela] (batch) at (0, 0) {%
    \texttt{export\_batch}
    \nodepart{two}
    \kwcolunas{\kwcol{id}{PK}\kwcol{received\_at}{TEXT}\kwcol{body\_bytes}{INTEGER}%
             \kwcol{point\_count}{INTEGER}}};

  \node[tabela, anchor=north west] (log) at (3.6, 2.6) {%
    \texttt{log\_record}
    \nodepart{two}
    \kwcolunas{\kwcol{id}{PK}\kwcol{batch\_id}{FK}\kwcol{received\_at}{TEXT}%
             \kwcol{service\_name}{TEXT}\kwcol{service\_instance\_id}{TEXT}%
             \kwcol{pipeline\_hop}{TEXT}\kwcol{scope\_name}{TEXT}%
             \kwcol{time\_unix\_nano}{INTEGER}\kwcol{observed\_time\_unix\_nano}{INTEGER}%
             \kwcol{body}{TEXT}\kwcol{event\_id}{INTEGER}%
             \kwdestaque{attributes}{TEXT (JSON)}%
             \kwcol{dedupe\_key}{TEXT, UNIQUE}}};

  \node[tabela] (raw) at (0, -3.2) {%
    \texttt{log\_record\_raw}
    \nodepart{two}
    \kwcolunas{\kwcol{id}{PK}\kwcol{batch\_id}{FK}\kwcol{received\_at}{TEXT}%
             \kwcol{reason}{TEXT}\kwcol{raw}{TEXT}}};

  \draw[fk] (log.west |- 0, 1.25) -| node[above, font=\scriptsize, pos=0.25] {\texttt{batch\_id}} (batch.north);
  \draw[fk] (raw.north) -- node[right, font=\scriptsize] {\texttt{batch\_id}} (batch.south);

  \node[nota, anchor=north west, text width=6.6cm] at ([yshift=-0.2cm]log.south west)
    {\texttt{attributes}: todos os atributos do registro, juntos, num
     documento JSON.\\
     \texttt{dedupe\_key}: chave de repetição; um registro repetido não é
     gravado de novo.};
\end{tikzpicture}

  \caption{O banco do backend próprio. Os atributos de cada registro ficam
    juntos num documento JSON; não há coluna por atributo nem tabela por
    fonte.}
  \label{fig:esquema-backend}
\end{figure}

Em seguida, o backend reduz de novo a granularidade do tempo à semana, para
que a regra valha também contra um cliente antigo ou de outra origem
(Seção~\ref{sec:mapeamento}); o Collector, como vem, não altera o tempo, e por
isso a regra fica no backend. É a única regra por campo aplicada no Collector
ou no backend. Com o tempo já reduzido, o backend calcula para cada registro
uma chave de repetição. Um registro com identificador de evento, como os que o
kw\_tele envia, é reconhecido pela instalação, pelo seu corpo e por esse
identificador; como o tempo deixou de distinguir um evento de outro, é o
identificador que faz esse papel.

Por fim, o backend grava o lote inteiro de uma vez, ignorando os registros
cuja chave já existe, e só então responde com sucesso; um reenvio, portanto,
não duplica nada. Se a gravação falha, ele responde com erro, e nada daquele
lote conta como recebido. Essa ordem é o que sustenta a cadeia de confirmação de ponta a ponta
(Figura~\ref{fig:confirmacao}). O kw\_tele só marca uma linha como enviada quando recebe uma
resposta de sucesso. O Collector só responde com sucesso depois que o backend
respondeu, porque não tem lote nem fila que o deixem responder antes. E o
backend só responde depois de gravar. Assim, um sucesso visto pelo kw\_tele
quer dizer que o lote foi gravado, com o que não coube no formato guardado à
parte.

Se o backend está parado, o Collector devolve erro, o kw\_tele não marca nada
e as linhas continuam pendentes; o próximo envio as manda de novo, e o
descarte de repetidos garante que um reenvio depois de uma falha parcial não
duplica registros. O padrão do exportador do Collector é ligar fila e novas
tentativas \citep{otel_collector_exporterhelper_v0_121_0}; elas foram
desligadas aqui para que a resposta de sucesso só venha depois da gravação.
Uma fila persistente no Collector, alternativa para uso em produção, não foi
construída.

\begin{figure}[!ht]
  \centering
  \resizebox{\textwidth}{!}{\begin{tikzpicture}[
    x=1cm, y=1cm,
    every node/.style={font=\footnotesize, align=center},
    ator/.style={draw, rounded corners=3pt, minimum height=0.7cm, minimum width=2.6cm,
                 inner sep=3pt, fill=white},
    dado/.style={ator, very thick, fill=black!12},
    vida/.style={densely dashed, thin},
    msg/.style={-{Latex[length=1.6mm]}, semithick},
    resp/.style={-{Latex[length=1.6mm]}, semithick, densely dotted},
    nota/.style={font=\scriptsize, inner sep=1.5pt},
    acao/.style={draw, fill=white, font=\scriptsize, inner sep=2pt, text width=2.4cm},
    quadro/.style={draw, rounded corners=2pt, inner sep=6pt},
    etiqueta/.style={draw, fill=white, font=\scriptsize\itshape, inner sep=2pt, anchor=north west},
  ]
  \foreach \n/\x/\t in {k/0/kw\_tele, c/4/Collector, b/8/backend próprio} {
    \node[ator] (\n) at (\x, 0) {\t};
    \draw[vida] (\n.south) -- (\x, -9.6);
  }
  \node[dado] (s) at (12, 0) {banco SQLite};
  \draw[vida] (s.south) -- (12, -9.6);

  \draw[msg] (0, -1.0) -- node[nota, above] {envia o lote} (4, -1.0);
  \draw[msg] (4, -1.5) -- node[nota, above] {repassa} (8, -1.5);
  \draw[msg] (8, -2.0) -- node[nota, above] {grava o lote} (12, -2.0);
  \draw[resp] (12, -2.5) -- node[nota, above] {gravado} (8, -2.5);
  \draw[resp] (8, -3.0) -- node[nota, above] {sucesso} (4, -3.0);
  \draw[resp] (4, -3.5) -- node[nota, above] {sucesso} (0, -3.5);
  \node[acao, anchor=north] at (0, -3.8) {marca o lote como enviado};

  \draw[msg] (0, -6.0) -- node[nota, above] {envia o lote} (4, -6.0);
  \draw[msg] (4, -6.5) -- node[nota, above] {repassa} (7.6, -6.5);
  \node[font=\footnotesize\bfseries] at (7.8, -6.5) {$\times$};
  \node[nota, right] at (8.1, -6.5) {sem resposta};
  \draw[resp] (4, -7.1) -- node[nota, above] {erro} (0, -7.1);
  \node[acao, anchor=north, text width=2.8cm] at (0, -7.4) {não marca; as linhas ficam pendentes e vão na próxima vez};
  \node[quadro, fit={(-1.6, -5.3) (13.4, -9.3)}] (alt) {};
  \node[etiqueta] at (alt.north west) {backend parado};

  \node[nota, anchor=west, font=\scriptsize\itshape, fill=white] at (1.5, -4.2) {a marca só vem depois do sucesso};
\end{tikzpicture}
}
  \caption{A confirmação. O kw\_tele só marca o lote como enviado depois que o
    sucesso volta do backend próprio, pelo Collector; com o backend parado, as
    linhas ficam pendentes.}
  \label{fig:confirmacao}
\end{figure}

\section{Consulta dos dados}
\label{sec:consulta}

O caminho descrito até aqui termina numa linha gravada pelo backend, e essas
linhas só servem se puderem ser lidas. O artigo do kw lista como alvos de
medição, por exemplo, em que hardware se compila e quem está sobrecarregado na
hierarquia de mantenedores (Seção~\ref{sec:dados-de-uso}); outros autores vão
querer outras coisas, e o backend não sabe de antemão que tipos de registro nem
que atributos vão chegar (Seção~\ref{sec:backend}). Por isso nenhuma consulta é
escrita para uma fonte: a pergunta chega como parâmetro. O backend próprio, o
mesmo que recebe os registros do Collector, oferece para isso três rotas, todas
só de leitura. Duas descrevem o que há nos dados, e a terceira agrupa e conta
(Figura~\ref{fig:consulta}).

\begin{figure}[!ht]
  \centering
  \resizebox{\textwidth}{!}{\begin{tikzpicture}[
    x=1cm, y=1cm,
    every node/.style={font=\footnotesize},
    pedido/.style={draw, rounded corners=3pt, inner sep=4pt, fill=black!6, align=center, text width=6.4cm},
    resp/.style={draw, inner sep=3pt, fill=white},
    dado/.style={draw, very thick, rounded corners=3pt, fill=black!12, inner sep=5pt, align=center},
    seta/.style={-{Latex[length=1.6mm]}, semithick},
    nota/.style={font=\scriptsize, align=center, inner sep=1.5pt},
  ]
  \node[pedido, anchor=north] (q1) at (0, 0) {\textbf{que tipos há?}\\[1pt]{\scriptsize\texttt{GET /analysis/schemas}}};
  \node[resp, anchor=north] (r1) at ([yshift=-0.5cm]q1.south) {%
    \setlength{\tabcolsep}{4pt}%
    \begin{tabular}{@{}lrr@{}}
      \scriptsize\textit{tipo} & registros & instalações \\
      \hline
      \texttt{build\_modules} & 112 & 9 \\
      \texttt{executed}       & 815 & 30 \\
    \end{tabular}};

  \node[pedido, anchor=north] (q2) at ([yshift=-0.7cm]r1.south) {\textbf{que atributos tem \texttt{executed}?}\\[1pt]
    {\scriptsize\texttt{GET /analysis/schema/executed/attributes}}};
  \node[resp, anchor=north] (r2) at ([yshift=-0.5cm]q2.south) {%
    \setlength{\tabcolsep}{4pt}%
    \begin{tabular}{@{}lr@{}}
      \scriptsize\textit{atributo} & registros \\
      \hline
      \texttt{kw.event.date}                       & 815 \\
      \texttt{kw.event.id}                         & 815 \\
      \texttt{kw.event.kind}                       & 815 \\
      \texttt{kw.executed.command\_label}          & 815 \\
      \texttt{kw.executed.elapsed\_time\_in\_secs} & 815 \\
      \texttt{kw.executed.status}                  & 815 \\
    \end{tabular}};

  \node[pedido, anchor=north, text width=7.0cm] (q3) at (8.6, 0) {\textbf{quantos módulos tem, em média, uma compilação}\\
    \textbf{que deu certo e uma que falhou?}\\[1pt]
    {\scriptsize\texttt{GET /analysis/query?schema=build\_modules}}\\[-1pt]
    {\scriptsize\texttt{\&join=executed\&by=kw.executed.status}}\\[-1pt]
    {\scriptsize\texttt{\&measure=kw.build\_modules.module\_count}}};
  \node[resp, anchor=north] (r3) at ([yshift=-0.5cm]q3.south) {%
    \setlength{\tabcolsep}{5pt}%
    \begin{tabular}{@{}lrrr@{}}
      estado & registros & instalações & módulos (média) \\
      \hline
      \texttt{failure} &   8 & 3 & 198 \\
      \texttt{success} & 104 & 9 & 212 \\
    \end{tabular}};
  \node[nota, anchor=north, text width=6.8cm] (j) at ([yshift=-0.15cm]r3.south)
    {o \texttt{join} junta os registros de tipos diferentes\\que vêm do mesmo evento};

  \node[draw, very thick, rounded corners=3pt, fill=white, inner sep=5pt, minimum width=16.4cm, align=center] (back) at (4.3, -10.0) {backend próprio};
  \node[dado, minimum width=5.0cm] (db) at (4.3, -11.3) {banco SQLite};
  \draw[seta] (db.north |- back.south) -- node[nota, right] {lê} (db.north);

  \draw[seta] (q1) -- (r1);
  \draw[seta] (q2) -- (r2);
  \draw[seta] (q3) -- (r3);
  \draw[seta] (back.north -| r2.south) -- node[nota, right] {responde} (r2.south);
  \draw[seta] (back.north -| r3.south) -- node[nota, right] {responde} (j.south -| r3.south);

\end{tikzpicture}
}
  \caption{As rotas de consulta sobre o banco do backend próprio. Valores
    ilustrativos.}
  \label{fig:consulta}
\end{figure}

A primeira rota devolve os tipos de registro que existem, isto é, os corpos,
que são os nomes das tabelas de extensão (Seção~\ref{sec:mapeamento}), cada
um com o número de registros e de instalações distintas. A segunda recebe um
desses tipos e devolve os atributos que aparecem nele, com quantos registros
têm cada um. É o ponto de partida quando não se sabe o que o kw grava: um tipo novo, como o da lista de módulos, aparece na primeira rota assim
que o primeiro registro dele é gravado, e os seus atributos, na segunda
(Seção~\ref{sec:extensibilidade}).

A terceira rota recebe um tipo e agrupa os seus registros pelos atributos
pedidos; se pedido, também resume uma medida numérica, como a duração de uma
execução. Os grupos podem ainda ser separados por semana, a granularidade com
que o tempo chega, e por instalação. A instalação aparece como um rótulo de ordem, e não com o seu identificador; o
risco desse rótulo está no Capítulo~\ref{cap:privacidade}. Antes
de agrupar, a consulta pode restringir os registros a um valor de atributo e a
um intervalo de semanas. Cada grupo traz o número de registros e o número de
instalações distintas. Os valores são agrupados como texto, e a medida só é
somada onde o valor é numérico, com os demais contados à parte, porque os atributos são guardados sem tipo.

Como cada fonte é um tipo separado, uma pergunta que cruza fontes precisa
juntar os registros. A terceira rota aceita outros tipos para juntar ao
pedido, ligando os registros do mesmo evento pela instalação e pelo identificador do
evento, que o cliente envia em todo registro e que o usuário não pode excluir
(Seção~\ref{sec:mapeamento}); um registro sem instalação ou sem esse
identificador não se junta a nenhum outro.
Com os dados de demonstração, a contagem de módulos vem dos registros da lista
de módulos e o estado da execução vem dos de execução. Juntando os dois tipos,
uma pergunta possível é se as compilações que falham têm mais módulos que as
bem-sucedidas; outra, sem junção, é quantas compilações terminaram em cada
estado a cada semana.

A consulta só devolve agregados: nenhuma rota devolve um registro individual
nem o identificador real de uma instalação, e os registros brutos ficam com
quem opera o backend. Ela também não traz painel nem gráfico. A resposta é um documento JSON, para
ser lida por um programa ou levada a uma ferramenta de análise, que também
poderia ler o mesmo banco SQLite; na leitura deste trabalho, é o mínimo para
responder a uma pergunta sem escrever código para ela. A falta de supressão de
grupos pequenos e de autenticação nas rotas são limites tratados no
Capítulo~\ref{cap:privacidade}, e a consulta sobre dados de várias
instalações é mostrada no Capítulo~\ref{cap:validacao}. A mesma consulta serve
a qualquer fonte que chegue pelo caminho (Seção~\ref{sec:extensibilidade}).

\section{Extensibilidade do sistema}
\label{sec:extensibilidade}

O sistema é extensível porque o conhecimento sobre o dado mora num ponto só, o
esquema do kw, e o resto do caminho o descobre em vez de o trazer escrito. A
Tabela~\ref{tab:extensao} resume o argumento, componente a componente, para uma
fonte nova. O kw é o único lugar que muda: a fonte nova é uma tabela ligada ao
evento e o código do kw que a preenche. O kw\_tele não precisa ser alterado,
porque descobre no esquema do banco as tabelas que estendem o evento e
transforma cada coluna num atributo (Seção~\ref{sec:mapeamento}). O protocolo
também não, porque ele define os atributos do registro de \textit{log} como
pares de chave e valor, sem fixar as chaves
\citep{otel_common_spec_v1_63_0,otlp_spec_1_11_1}, e um tipo novo é apenas
outro valor no corpo do registro. O Collector não olha os atributos. A exceção é o tempo: uma coluna de tempo própria da extensão sai como foi
gravada, porque a redução da granularidade só vale para as colunas do evento
(Seção~\ref{sec:mapeamento}).

\begin{table}[!ht]
  \centering
  \footnotesize
  \setlength{\tabcolsep}{4pt}
  \begin{threeparttable}
    \caption{O que muda em cada componente quando entra uma fonte nova, como a
      lista de módulos da compilação. Só o kw muda; a exceção é uma coluna de tempo na tabela nova, que não tem
      a granularidade reduzida no cliente.}
    \label{tab:extensao}
    \begin{tabular}{@{}>{\raggedright\arraybackslash}p{2.4cm}
                      >{\raggedright\arraybackslash}p{4.6cm}
                      >{\raggedright\arraybackslash}p{7.2cm}@{}}
      \toprule
      Componente & O que muda para uma fonte nova & Por quê \\
      \midrule
      kw & uma tabela cuja chave é o identificador do evento, e o código que a preenche
         & é onde o dado nasce \\
      \addlinespace[2pt]
      kw\_tele & nada
         & descobre no esquema as tabelas que estendem o evento \\
      \addlinespace[2pt]
      protocolo & nada
         & o registro de log carrega atributos livres, e o tipo é um valor do corpo \\
      \addlinespace[2pt]
      Collector & nada
         & não olha os atributos \\
      \addlinespace[2pt]
      backend próprio & nada
         & guarda os atributos juntos, sem uma coluna por atributo \\
      \addlinespace[2pt]
      consulta & nada
         & descobre nos dados os tipos e os atributos, e recebe o nome do
           atributo como parâmetro \\
      \bottomrule
    \end{tabular}
    \begin{tablenotes}
      \footnotesize
      \item Justificativas na Seção~\ref{sec:extensibilidade}.
    \end{tablenotes}
  \end{threeparttable}
\end{table}


No backend próprio, os atributos de cada registro são guardados juntos, num
documento JSON, e não numa coluna para cada um; a tabela de registros não
depende, assim, de que atributos chegam nem de como se chamam
(Figura~\ref{fig:esquema-backend}). A consulta descobre nos próprios dados os
tipos de registro e os atributos que existem, recebendo o nome do atributo como
parâmetro (Seção~\ref{sec:consulta}). Nenhum componente depois do kw conhece
fonte alguma, e por isso não há nada a registrar quando uma fonte aparece. A
privacidade acompanha o dado novo do mesmo modo: a exclusão de campos é uma
lista de chaves de atributo, e uma chave listada vale para o atributo que a
tiver, exista ele ou não no momento em que o usuário a configurou
(Capítulo~\ref{cap:privacidade}). Esse é o critério de extensibilidade da
Seção~\ref{sec:criterios} visto pelo lado do design, e é o que o sexto objetivo
deste trabalho pede (Seção~\ref{sec:objetivos}).

Um exemplo torna o argumento concreto: a lista de módulos de uma compilação,
isto é, os módulos que a configuração do kernel marca para serem compilados à
parte. Ela foi acrescentada ao kw como exemplo do argumento, e não por ser um
dado escolhido para coleta. No kw, e só nele, ela é uma tabela ligada ao evento da compilação ou da
instalação e o código que a preenche (Seção~\ref{sec:banco-local}). Uma
compilação ou uma instalação passa a gerar dois registros com o mesmo
identificador de evento, um por tabela, e a chave de repetidos do backend, que
inclui o corpo, os mantém separados. A prova de que o exemplo passa pelo
caminho, sem mudança no kw\_tele nem no backend, está no
Capítulo~\ref{cap:validacao}.
